\chapter{ระบบตรวจจับและสั่งการไร้สัมผัสด้วย AR MediaPipe Hands}

การควบคุมการทดลองวิทยาศาสตร์เสมือนจริงโดยไม่ใช้เมาส์หรือคีย์บอร์ด ช่วยเสริมสร้างทักษะการลงมือปฏิบัติจริงเชิงประจักษ์ (Kinesthetic Learning) ผ่านการวิเคราะห์โครงข่ายประสาทเทียมของมือมนุษย์แบบเรียลไทม์

\section{โครงสร้างเรขาคณิตของมือมนุษย์ 21 จุดข้อต่อ}

โมเดลคอมพิวเตอร์วิทัศน์ MediaPipe Hands สกัดจุดพิกัดสามมิติ $(x, y, z)$ ของข้อนิ้วจำนวน 21 จุด โดยมีจุดอ้างอิงสำคัญ ได้แก่:
\begin{itemize}
    \item \textbf{WRIST (จุดที่ 0):} ข้อมือส่วนล่างสำหรับกำหนดเวกเตอร์การเคลื่อนที่ของมือ
    \item \textbf{THUMB\_TIP (จุดที่ 4):} ปลายนิ้วหัวแม่มือ
    \item \textbf{INDEX\_FINGER\_TIP (จุดที่ 8):} ปลายนิ้วชี้สำหรับใช้เป็นเคอร์เซอร์ระบุตำแหน่ง
    \item \textbf{MIDDLE\_FINGER\_TIP (จุดที่ 12):} ปลายนิ้วกลางสำหรับประเมินการกำมือ
\end{itemize}

\section{อัลกอริทึมการคำนวณเวกเตอร์ท่าทางมือ}

การจำแนกท่าทาง (Gesture Classification) ใช้การคำนวณระยะทางแบบยูคลิด (Euclidean Distance) ระหว่างปลายนิ้วชี้และปลายนิ้วหัวแม่มือ:

\begin{equation}
d_{\text{pinch}} = \sqrt{(x_8 - x_4)^2 + (y_8 - y_4)^2 + (z_8 - z_4)^2}
\end{equation}

เมื่อ $d_{\text{pinch}} < \delta_{\text{threshold}}$ (กำหนดค่าที่ $0.05$ ในพิกัดนอร์แมลไลซ์) ระบบจะส่งสัญญาณสถานะ \texttt{PINCH\_ACTIVE} เพื่อจำลองการหยิบจับสายไฟ ขยับลูกตุ้ม หรือหมุนปุ่มปรับค่าบนแผงควบคุม

\section{การผสานระบบเสียงอคูสติกสดผ่าน Web Audio API}

เมื่อมือของผู้ใช้สัมผัสหรือปล่อยอุปกรณ์ ระบบจะสังเคราะห์คลื่นเสียงความถี่เฉพาะผ่าน Web Audio API โดยไม่ต้องโหลดไฟล์เสียงภายนอก:

\begin{equation}
f_{\text{click}} = 880\,\text{Hz}, \quad \tau_{\text{decay}} = 0.05\,\text{s}
\end{equation}

ทำให้ผู้เรียนได้รับข้อมูลป้อนกลับแบบสองทาง (Haptic Audio Feedback) ที่ตอบสนองในระดับมิลลิวินาที
